\title{Large Language Models Can Automatically Engineer Features for Few-Shot Tabular Learning}

\begin{abstract}
Large Language Models (LLMs) demonstrate exceptional performance in complex and novel reasoning scenarios, presenting significant promise for tabular learning—a foundational technique in numerous practical domains. We introduce a new in-context learning methodology, \textsf{FeatLLM}, which leverages LLMs as automated feature engineers to construct data representations tailored for optimized tabular prediction tasks. The transformed features are subsequently utilized to estimate class probabilities with a basic downstream predictive model (e.g., linear regression), delivering robust results even in low-data (few-shot) contexts. At prediction time, \textsf{FeatLLM} relies exclusively on this lightweight predictor operating over the engineered features, in contrast with LLM-dependent methods that require runtime LLM queries for each instance. This design circumvents continuous LLM querying during inference, merely necessitating API-level access during training, and thus bypasses issues such as prompt length constraints. Empirical evaluations conducted across diverse tabular benchmarks reveal that \textsf{FeatLLM} produces concise, high-utility rules and, on average, achieves a 10\% performance boost relative to contemporaries such as TabLLM and STUNT.
\end{abstract}

\section{Introduction}
Large Language Models (LLMs) have, in recent times, displayed formidable competence in generating appropriate responses for tasks unobserved during training, all without necessitating specific fine-tuning for those tasks~\cite{creswell2022selection,imani2023mathprompter,taylor2022galactica}. Their pretraining on voluminous text datasets equips them with rich prior knowledge, enabling them to, in many cases, effectively replace specialized domain expertise~\cite{Genomes2021,singhal2023large,wilcox2003role}. This quality has positioned LLMs as critical enablers of automation in a wide array of applications—from conversational analysis~\cite{finch2023leveraging} to automated code repair~\cite{fan2023automated}. Notably, the inclusion of task-specific descriptions and a handful of illustrative examples within prompts allows LLMs to understand task context and focus on salient features~\cite{brown2020language,zhao2023survey}. These abilities illustrate that LLMs possess analytical acumen akin to skilled feature engineers.

Recently, efforts to harness LLMs’ prior knowledge and reasoning power have been directed towards tabular learning. For example, prior domain insight—such as age being a risk factor in disease prognosis—can significantly aid tasks like medical prediction. To capture these benefits, recent studies construct natural language descriptions of both tasks and features, serialize the tabular data, and introduce this information to LLMs~\cite{dinh2022lift,wang2023anypredict}. These pipelines often employ either in-context learning frameworks~\cite{nam2023semi} or leverage parameter-efficient tuning strategies~\cite{dinh2022lift,hegselmann2023tabllm}. In many low-resource situations, such injection of LLM-provided prior knowledge has consistently surpassed traditional tabular learning methods~\cite{hegselmann2023tabllm}.

Existing methods for leveraging LLMs in tabular learning encounter notable shortcomings. For direct, end-to-end prediction tasks, every data instance necessitates a separate LLM query, resulting in substantial computational costs. Furthermore, achieving state-of-the-art predictive performance typically demands LLM fine-tuning, which is impractical for those models whose training infrastructure or tuning APIs remain unavailable—an issue exacerbated as recent high-performing LLMs increasingly restrict users to inference-only APIs~\cite{anil2023palm,openai2023gpt4}. Another limitation is poor scalability to longer prompt lengths~\cite{liu2023lost}; as the feature dimensionality of tabular datasets grows, the tokenized representations often surpass practical limits or suffer from noticeable drops in predictive quality. Collectively, these restrictions present significant obstacles to integrating LLM-based tabular models in operational contexts.

To address these issues, our work departs from conventional end-to-end LLM usage and instead repurposes LLMs primarily for feature engineering, thereby exploiting their intrinsic domain knowledge more effectively. We propose an innovative in-context learning strategy called \textsf{FeatLLM}, with the objective of uncovering the ``criteria" or rationale the LLM uses when making predictions. For example, when diagnosing a disease, our framework prompts the LLM to articulate explicit rules—identifying which feature patterns are indicative of the outcome. These extracted rules then inform the synthesis of new features that can supplant the original attributes during downstream modeling. This paradigm reframes the contribution of the LLM towards creating derived features, subsequently enabling downstream models to become considerably simpler; inference can then proceed rapidly, without requiring a complex predictive pipeline. Utilizing the in-context learning capacity of LLMs, the process can forego traditional training, instead leveraging a small set of exemplar prompts. Additionally, to handle challenges posed by large input prompts, we integrate ensemble learning methods such as feature bagging~\cite{breiman2001random}, helping to maintain scalability.

\textsf{FeatLLM} decomposes the process of generating features for novel tasks into two distinct subtasks: (1) grasping the specific problem and deducing how features interact with targets, and (2) utilizing the insights from this understanding, together with limited labeled instances, to formulate critical rules that separate different categories. By organizing the reasoning in this manner, LLMs can better identify and exclude irrelevant features when developing extraction rules. The established rules are then applied to the dataset (after conversion into executable program code), transforming the original data into binary indicators that specify whether individual samples satisfy particular rules. Subsequently, a simple predictive model is trained to compute class probabilities based on these engineered binary features. Reliability is further enhanced through an ensemble methodology that iteratively regenerates features using feature bagging. Properly selecting the number of features and samples included during bagging helps address concerns about overly large prompt sizes. Free from the requirement to train models and imposing minimal inference costs, \textsf{FeatLLM} enables the practical adoption of LLMs across a broad array of applied problems.

Our assessment of \textsf{FeatLLM} covers 13 tabular datasets under limited sample conditions, demonstrating its effectiveness and resilience. The results confirm that LLMs adeptly draw upon both prior and empirical data knowledge to engineer high-quality features. The framework consistently surpasses existing state-of-the-art few-shot learning methods across multiple test scenarios. Code is accessible at the anonymized GitHub repository: \texttt{https://github.com/Sungwon-Han/FeatLLM}.

\section{Related Work}

\subsection{Few-Shot Learning with Tabular Data} Developments in few-shot learning have made it possible to derive representations that are broadly applicable, even when annotation is limited~\cite{chen2018closer}. While initial work was centered around visual data, few-shot learning techniques have exhibited considerable generalizability in other formats as well, such as text, audio, and, more recently, tabular data~\cite{majumder2022few,nam2023stunt,schick2021exploiting}. The challenge of learning from just a handful of labeled instances increases susceptibility to incidental patterns~\cite{bartlett2021deep}. Consequently, the field has gravitated toward supplementing model training with extra information sources or integrating domain knowledge to introduce appropriate inductive biases. Representative strategies include harnessing large unlabeled datasets supplemented with targeted augmentations for semi-supervised objectives~\cite{ucar2021subtab,yoon2020vime}, or obtaining superior initializations by employing unsupervised pre-training that is not task-specific~\cite{somepalli2021saint}. These unsupervised methods frequently rely on objectives like masking or introducing noise to portions of the data, then enforcing reconstruction~\cite{arik2021tabnet,majmundar2022met}, or they exploit data augmentation strategies for contrastive learning~\cite{bahri2022scarf,chen2020simple}. Still, the diversity intrinsic to tabular datasets complicates the design of universally applicable augmentations, often resulting in limited impact for these methods when label scarcity is severe~\cite{nam2023stunt}. 

More recent research has shifted focus to cross-dataset learning, wherein models are trained on a diverse set of tabular datasets unrelated to the task at hand, before deploying transfer learning to target problems~\cite{zhang2023generative,levin2022transfer,zhu2023xtab}. One approach, TransTab~\cite{wang2022transtab}, leverages transformer architectures to extract semantic connections by conditioning on column names in addition to table contents. The information gleaned from a variety of column types enables both zero-shot and few-shot generalization to previously unseen tabular tasks. Alternatively, TabPFN~\cite{hollmann2022tabpfn} constructs synthetic data spanning various distributions for the purpose of pre-training a Bayesian neural network, after which predictions are generated by jointly presenting the model with both the training and testing records from the target task—without requiring additional fine-tuning. The cross-dataset prior captured during pre-training has led to consistent improvements over conventional decision tree baselines, particularly under few-shot regimes.

\subsection{Language-Interfaced Tabular Learning} Integrating large language models (LLMs) presents another promising path for incorporating prior knowledge into tabular learning~\cite{anil2023palm,openai2023gpt4}. Owing to their training on massive, heterogeneous text sources, LLMs are capable of exhibiting transfer to new, unseen tasks and have proved beneficial across different disciplines~\cite{brown2020language}. Recent efforts have aimed to exploit the semantic and commonsense knowledge stored in LLMs for analyzing tabular data, often by converting structured tables into textual representations that are suitable for prompt-based inference~\cite{dinh2022lift,hegselmann2023tabllm,wang2023anypredict}. Within these prompts, inclusion of not only the tabular entries but also explanatory descriptions for features and tasks enables the model to infer complex task-feature relationships needed for prediction. Approaches diverge into parameter-efficient fine-tuning schemes such as LoRA~\cite{hu2021lora} and IA3~\cite{liu2022few}, and methods that utilize in-context learning, wherein prompt templates are augmented with a small set of demonstration examples to enable task adaptation without any explicit gradient-based optimization~\cite{dinh2022lift,hegselmann2023tabllm,nam2023semi,slack2023tablet}.

Although the previously mentioned approaches have demonstrated success in advancing few-shot tabular learning, their reliance on running inference through the LLM each time remains a barrier for deployment in industrial contexts. This work proposes a shift away from the standard paradigm, which uses the LLM in a black-box, end-to-end manner to generate predictions for individual samples. Rather, our method leverages the LLM to explicitly deduce the rules or criteria that govern each class. \textsf{FeatLLM} translates these LLM-generated rules into executable code to engineer new feature sets, thereby facilitating predictions independent of the LLM. By utilizing in-context learning, \textsf{FeatLLM} harnesses both background knowledge and information derived from limited labeled samples to extract rules.

\section{Methods}
\textsf{FeatLLM} is built to identify and extract class-defining “rules” from a handful of example inputs, as illustrated in Figure~\ref{fig:main_model}. These rules, inferred via the LLM, are subsequently converted into programmatic logic that yields binary feature values for new instances, flagging whether a particular rule applies. The resulting collection of binary features undergoes a weighted dot product with non-negative, learnable parameters to calculate the probability of each class assignment. Training this model allows for learning the significance of each derived feature in predicting class likelihood (see Section~\ref{Sec:training}). To enhance stability and robustness, the procedure is performed repeatedly with bagging and the outputs are then aggregated through ensembling. Further elaboration on the problem formulation and the details of each step follow below.

\vspace*{-0.12in}
\paragraph{Problem formulation.} We examine a tabular dataset comprising $N$ labeled instances, denoted as $\mathcal{D} = \{(\mathbf{x}^i, \mathbf{y}^i)\}_{i=1}^{N}$. Each instance is represented by a $d$-dimensional feature vector $\mathbf{x}^i$, and the set of input features, $F = \{f_j\}_{j=1}^{d}$, is annotated with their corresponding natural language names and detailed definitions supplied externally. Within the context of a classification task, each label $\mathbf{y}^i$ indicates membership in one of several predefined categories. In the $k$-shot learning context, we randomly select $k$ labeled examples, where $k < N$, to serve as the training data.

\subsection{Prompt Design for Extracting Rules}
\label{Sec:prompt_design}
To facilitate more precise and logical rule extraction by the LLM, we craft prompts to simulate the reasoning process an expert would follow when confronted with tabular data analysis. Our prompt structure integrates three principal elements (illustrated in Fig.~\ref{fig:prompt_for_rule} and detailed in Appendix~\ref{sec:example_rule}):

\vspace*{-0.12in}
\paragraph{Basic information description.} This initial section offers foundational details crucial for addressing the task at hand (refer to the orange-highlighted segment in Fig.~\ref{fig:prompt_for_rule}). It outlines the problem statement along with target label specifics, definitions for each input feature, and showcases example instances. The objective is presented in a question format (e.g., ``Does the information regarding this patient's myocardial infarction complications suggest chronic heart failure? Yes or no?". Supplementary task descriptions are available in Appendix~\ref{sec:task_desc}). For every feature, we specify its data type and, when applicable, provide a concise definition; categorical features may also have representative category values listed. Sample demonstrations are created by converting select training instances into text sequences, appending the correct labels. The serialization adheres to conventions established by prior literature~\cite{dinh2022lift,hegselmann2023tabllm}:

\begin{align}
&\text{Serialize}(\mathbf{x}^i,\mathbf{y}^i, F) = \nonumber \\ 
&``\ f_1\ \text{is}\ \mathbf{x}^i_1.\ ... \ f_d\  \text{is}\  \mathbf{x}^i_d.\ \ \text{Answer:}\  \mathbf{y}^i\ ”
\end{align}

\vspace*{-0.12in}
\paragraph{Reasoning instruction.} Our approach seeks to improve the reasoning capabilities of LLMs by supplying them with explicit reasoning instructions (illustrated by the blue text in Fig.~\ref{fig:prompt_for_rule}). This begins with an introductory statement, paralleling the chain-of-thought paradigm~\cite{wei2022chain}. We further decompose the LLM's reasoning into two sequential phases. First, the LLM is prompted to determine causal associations or inclinations between input features and the task description, drawing upon domain knowledge or intuition instead of relying on specific examples. This enables the LLM to harness its background knowledge relevant to the problem. In the subsequent phase, the LLM integrates the provided example demonstrations alongside its prior inferences from the first step to construct classification rules. By structuring the inference in two phases, we reduce the likelihood of the model focusing on irrelevant or coincidentally correlated columns and facilitate attention toward crucial features. To balance prompt length and prevent both under- and overfitting due to rule quantity, we standardize the number of rules generated per class (set at 10)\footnote{For additional details regarding rule quantity, refer to Section~\ref{sec:analysis}.}. Additionally, during rule construction, the LLM references the datatype information provided in the task description to ensure compatibility between feature types and rule expressions.

\vspace*{-0.12in}
\paragraph{Response instruction.} To support efficient parsing and downstream application of the generated rules, we provide the LLM with clear response structuring instructions (denoted by the yellow text in Fig.~\ref{fig:prompt_for_rule}). The prompt spells out the formatting requirements for responses corresponding to each class label (i.e., the expected answer format) and specifies the preferred structure for individual rules (i.e., the rule template). These directives help the LLM to select the correct rule format based on feature type. If further structure is not provided, the LLM is allowed to combine rules with logical connectives such as AND or OR. An exemplar output is shown in Appendix~\ref{sec:outcome-rule}.

\subsection{Inferring Class Likelihood via Rules}
\label{Sec:training}

\paragraph{Parsing rules for feature generation.} The rules established in the preceding phase are employed to derive novel binary features. For each target class, these binary features signal whether a data point adheres to the specific set of rules for that class (assigned as '0' or '1'). Such features are subsequently useful for modeling the probability that a sample pertains to a given class. Nevertheless, because the LLM expresses these rules in natural language, automated feature extraction requires a systematic parsing mechanism. While formatting guidelines are imposed during the rule elicitation stage to facilitate subsequent parsing, LLM outputs often still exhibit syntactic noise and unexpected textual modifications~\cite{kovavc2023large}. For example, the LLM may inconsistently replace symbols like ``$>$” or ``$<$” with their verbal equivalents such as ``greater than'' or ``less than'', thus complicating the parsing process.

Rather than engineering elaborate parsing logic to handle noisy textual output, we directly harness the capabilities of the LLM. The LLM’s semantic comprehension allows it to recognize the intended meaning, even in the presence of stylistic or syntactic shifts, and its pretraining on code datasets enables code generation in response to detailed instructions. To utilize these advantages, we craft prompts to include the function definition, explicit input and output formats, and the set of inferred rules, submitting these to the LLM for interpretation. Example prompts and their resulting code are provided in Figures~\ref{fig:prompt_for_parsing} and~\ref{sec:example_parsing}-\ref{sec:example_parse_outcome} in the Appendix. The LLM’s output is subsequently executed in Python using the \texttt{exec()} statement, thereby implementing the conversion of input data based on the parsed rules.

\vspace*{-0.12in}
\paragraph{Inferring class likelihood.} When new binary features are obtained from rules for each class, one straightforward approach for assessing how likely a sample belongs to a particular class is to tally the count of satisfied rules corresponding to each class (i.e., computing the sum of that class's binary features). Yet, this does not account for the fact that individual rules and features may influence the outcome differently. Therefore, their relevance must be estimated from the training data. To determine their impact, we employ a standard machine learning (ML) model on the binary features specific to each class. For instance, a linear model without an intercept can be employed. Suppose $\mathbf{z}_k^i \in \{0, 1\}^R$ denotes the set of binary features extracted from the $i$-th sample for class $k$, where $R$ indicates the count of rules per class and $c$ is the total number of classes. The class probability vector $\mathbf{p}^i$ for the sample is then determined as:

\begin{align}
\text{logit}_k^i &= \max(\mathbf{w}_k, 0) \cdot \mathbf{z}_k^i, \nonumber \\
\mathbf{p}^i &= \text{Softmax}({[\text{logit}_{1}^i, ..., \text{logit}_{c}^i]}). \label{eq:infer_prob}
\end{align}

In this context, $\mathbf{w}_k$ are trainable parameters within the linear model that denote the relative importance of each feature. By ensuring the weight vector only covers the non-negative domain, the LLM-derived features are prevented from adversely affecting a class's likelihood during the learning process. The resulting logit scores are transformed to a distribution over classes using the softmax function. Model training involves minimizing the cross-entropy loss based on the few-shot labeled data to update the weight vectors $\mathbf{w}_k$.

\vspace*{-0.12in}
\paragraph{Ensembling with bagging.} To further enhance robustness and predictive performance, the aforementioned procedure is run repeatedly to generate multiple inference models, which are then aggregated for the final decision through an ensemble approach. Let $\{\mathbf{w}[t]\}_{t=1}^T$ be the ensemble of weight vectors obtained from $T$ runs. For each inference, class probabilities are calculated using each $\mathbf{w}[t]$ (per Eq.~\ref{eq:infer_prob}), and the results are averaged across the ensemble to arrive at the final prediction.

To embed stochasticity into the rule creation process for ensembles, we set a relatively high temperature during LLM inference. In addition, we shuffle the sequence of few-shot examples in light of existing findings that the ordering can affect LLM responses~\cite{lu2022fantastically}\footnote{Analysis on the rule diversity can be found in Appendix~\ref{sec:diversity}}. To leverage predictive diversity for greater robustness, we apply bagging, randomly choosing subsets of features or samples in each trial—a strategy that is particularly advantageous when handling datasets with a substantial number of features or data points. Our proposed ensemble strategy offers two primary benefits. First, even when the LLM produces faulty rules during some trials, accurate rules emerging from other trials serve to mitigate these errors, thus bolstering the overall resilience of the system. This phenomenon is closely related to the self-consistency property of LLMs~\cite{wang2022self}; rules that repeatedly appear across independent trials have a higher probability of being correct. Second, our ensemble approach helps overcome the constraints imposed by limited prompt length. For instances where tabular datasets are too large to fit into a single prompt, the use of bagging effectively reduces input size while maintaining robust performance. Since a single set of rules may not represent the full relationships among all features, constructing an ensemble of several weak learners across trials minimizes the drawbacks of partial feature or instance inclusion within any single trial.

\section{Experiments}
We assess \textsf{FeatLLM} across multiple tabular datasets, accompanied by analyses intended to shed light on the workings of our approach. Space limitations preclude inclusion of certain experiments here (including comparisons of different LLMs and qualitative insights); these are presented in the Appendix.

\subsection{Performance Evaluation}
\paragraph{Datasets.}
For evaluation, we selected 13 datasets spanning binary and multiclass classification tasks: (1) Adult~\cite{asuncion2007uci}, which entails predicting whether an individual's income exceeds \$50,000 per year; (2) Bank~\cite{moro2014data}, where the goal is to forecast if a bank customer will open a term deposit account; (3) Blood~\cite{yeh2009knowledge}, focused on predicting a donor’s likelihood to return for another donation; (4) Car~\cite{kadra2021well}, aimed at assessing car quality; (5) Communities~\cite{misc_communities_and_crime_183}, used to estimate crime rates in geographic regions; (6) Credit-g~\cite{kadra2021well}, which involves determining if an individual presents a favorable or unfavorable credit risk; (7) Diabetes\footnote{\texttt{kaggle.com/datasets/uciml/pima-indians-diabetes-database}}, predicting an individual’s diabetes status; (8) Heart\footnote{\texttt{kaggle.com/datasets/fedesoriano/heart-failure-prediction}}, focused on forecasting the presence of coronary artery disease; and (9) Myocardial~\cite{misc_myocardial_infarction_complications_579}, tasked with predicting chronic heart failure in individuals.

In addition, we incorporate both recently introduced and synthetic datasets that have received limited attention in existing literature and were excluded from the pre-training phase of LLMs: (10) Cultivars, which aims to forecast the grain yield for particular soybean cultivars\footnote{\texttt{archive.ics.uci.edu/dataset/913}}; (11) NHANES, designed to determine a person's age group based on their record\footnote{\texttt{archive.ics.uci.edu/dataset/887}}; (12) Sequence-type, which assigns sequences to one of four categories: arithmetic, geometric, Fibonacci, or Collatz; and (13) Solution-mix, which classifies whether a mixture comprised of four solutions possesses a concentration exceeding 0.5. The first two datasets were donated to the UCI Machine Learning Repository post-September 2023, making certain they do not overlap with the data used in pre-training the LLM API (gpt-3.5-turbo-0613) employed by our approach. The latter two datasets are synthetic, thereby precluding the possibility of memorization by the LLM API and allowing fair assessment.

The datasets represent a broad spectrum in terms of size and difficulty, as outlined in Table~\ref{Tab:basic-desc}. Each offers clear attribute naming and concise descriptions. Datasets such as `Communities' and `Myocardial' include upwards of 100 features, which introduces additional complexity due to the LLM’s restricted context window.

\vspace*{-0.12in}
\paragraph{Baselines.}
We benchmark our results against nine existing approaches. Three classic tabular learning methods are considered: (1) Logistic Regression (LogReg), (2) XGBoost~\cite{chen2016xgboost}, and (3) RandomForest~\cite{ho1995random}. The following two baselines are tailored for scenarios with access only to unlabeled data: (4) SCARF~\cite{bahri2022scarf} and (5) STUNT~\cite{nam2023stunt}, although sourcing large unlabeled datasets may present practical difficulties. Additionally, (6) TabPFN~\cite{hollmann2022tabpfn} is evaluated; this method utilizes synthetically generated datasets with varying data distributions for model pretraining. The last three baselines leverage LLM-based techniques, namely (7) In-context learning (In-context)~\cite{wei2022emergent}, (8) TABLET~\cite{slack2023tablet}, and (9) TabLLM~\cite{hegselmann2023tabllm}. With In-context learning, training examples are incorporated into the prompt for inference without model parameter adjustments. TABLET augments the prompt with supplementary guidance derived from an external classifier using rule-sets and prototypes to improve inferential accuracy. Lastly, TabLLM applies a parameter-efficient tuning method, such as IA3~\cite{liu2022few}, to update the LLM with limited examples.

\vspace*{-0.12in}
\paragraph{Implementation details.}
\textsf{FeatLLM} adopts GPT-3.5 as its primary LLM, though it is purposefully constructed to remain independent of any specific LLM (results using various backbones can be found in Appendix~\ref{sec:backbones}). For inference, the temperature parameter is maintained at 0.5 while the top-p setting stays at the API’s standard value of 1. The ensemble size is fixed at 20, and the rule extraction process is performed with 10 rules. Figure~\ref{fig:hyperparameter2} and Appendix~\ref{sec:hyper-parameter} present a deeper analysis of how different hyper-parameters affect performance.

The linear model is optimized with Adam, set to a learning rate of 0.01, and is trained for 200 epochs. To select the best epoch, we rely on $k$-fold cross-validation. For baseline reproduction, in-context learning uses GPT-3.5, while TabLLM relies on T0~\cite{sanh2022multitask} according to its original configuration. Note that TabLLM’s use of LLMs is limited to those for which public checkpoints are provided, so GPT-3.5 is unavailable in these experiments. For classical machine learning algorithms like LogReg, XGBoost, and RandomForest, parameters are chosen via Grid-search in combination with $k$-fold cross-validation. The choice of $k$ (either 2 or 4) ensures every class is present in the training split. Additional implementation specifics can be found in Appendix~\ref{sec:baselines}.

\vspace*{-0.12in}
\paragraph{Results.}
Tables~\ref{tab:main1} and \ref{tab:main2} display comparative results on various datasets. Each experiment is conducted three times using different random splits for training and testing, and both the mean and standard deviation of the resulting AUC scores are reported. \textsf{FeatLLM} is either the best or second-best performer among all methods considered. Notably, our approach achieves results on par with standard tabular baselines requiring 32 shots, even when only 4 shots are available (refer to Fig.~\ref{fig:compares}). While approaches like STUNT and TabPFN exhibit strong gains on specific datasets, their aggregate performance improvements over traditional machine learning methods are limited. Additionally, LLM-based models—including in-context learning, TABLET, and TabLLM—were unable to process tabular data with many features. In contrast, our methodology leverages bagging and ensemble methods to handle datasets with numerous features effectively, yielding robust performance. It is also worth mentioning that adapting our bagging and ensemble strategy to other LLM-based methods is theoretically possible. However, this would double the computational cost per data point, substantially increasing the inference workload and rendering such adaptations less practical. 

\subsection{Analysis \& Discussion}
\label{sec:analysis}
\paragraph{Ablation study.} We conducted a series of ablation experiments to examine how key components influence the system's performance, focusing specifically on: (1) \textsf{-Tuning}: eliminating the weight tuning procedure in the linear model, instead directly summing the binary feature outputs to obtain logits; (2) \textsf{-Ensemble}: removing the ensemble step; (3) \textsf{-Description}: excluding the feature description; and (4) \textsf{-Reasoning}: removing the Step-1 phase of reasoning instruction during rule generation.

Table~\ref{tab:ablation} details the variations in AUC resulting from these ablations, averaged over all datasets. In every case, omission of a component led to diminished performance. Notably, the impact of weight tuning becomes more significant as the number of shots increases, implying that rule importance can be assessed more accurately with larger datasets and that tuning accordingly provides greater benefit. Conversely, omitting feature descriptions and reasoning instructions is more detrimental at lower shot counts, indicating that effective harnessing of LLM prior knowledge is particularly advantageous when limited training data is available. Finally, incorporating our ensemble method yields consistent and significant improvements across all experimental configurations.

\vspace*{-0.12in}
\paragraph{Training \& inference time comparison.} We evaluate and contrast the duration required for both training and inference across multiple approaches. All experiments utilize the Adult dataset with a 16-shot training set. By default, a single A100 GPU is allocated, while TabLLM leverages four A100 GPUs to facilitate model parallelism. Table~\ref{tab:cost} summarizes the comparative runtimes. Importantly, \textsf{FeatLLM} achieves relatively low inference latency, on par with standard baselines. On the other hand, LLM-driven methods including In-context learning, TABLET, and TabLLM are associated with substantially longer inference durations. For training time, \textsf{FeatLLM} remains competitive relative to other LLM-oriented techniques, as it requires just 30 API queries, posing little budgetary concern and enabling parallel execution.

\vspace*{-0.12in}
\paragraph{Quality of features generated by \textsf{FeatLLM}.} 
We implement a straightforward experiment to evaluate how informative the rule-based features produced by \textsf{FeatLLM} are for practical machine learning tasks. In particular, we calculate the average AUROC for the features generated and their association with target labels, then determine the proportion of features on each dataset exhibiting an AUROC above 0.5. Table~\ref{tab:quality} presents these outcomes. The data indicate that a significant portion of generated rules are pertinent to the target variable, thus contributing valuable information to the prediction task (see “Before tuning” in the table). Furthermore, while fitting the linear model, rules lacking strong connections to the data—identified by weights below a specified threshold—are naturally omitted (as reflected in “After tuning”). Here, the threshold for weights was chosen as 0.1 based on the overall weight histogram.

\vspace*{-0.12in}
\paragraph{Effect of adding spurious correlations.} We examined how well different approaches mitigate the influence of irrelevant spurious correlations under conditions of limited labeled data. For this analysis, we selected the Heart dataset and artificially introduced noise by appending randomly chosen columns from the Adult dataset, which bear no connection to the predictive task of the Heart dataset. Subsequently, we evaluated how model accuracy is affected as the number of these extraneous columns grows. To maintain an equitable comparison, each method received only 8 labeled examples, with no access to additional unlabeled data, thereby omitting techniques such as SCARF or STUNT from consideration.

Figure~\ref{fig:spurious} illustrates how performance shifts as spurious correlations are introduced. Notably, \textsf{FeatLLM} exhibited the smallest decrease in accuracy among all methods, suggesting that it is particularly resilient to this form of noise. This robustness likely stems from the fact that our approach compels attention to the relevance of each feature for the specific task by leveraging prior knowledge during rule derivation, thereby diminishing the likelihood that rules rely on unrelated columns. Supporting this, we found that only about 1-2\% of derived rules referenced the artificially added irrelevant features. Nevertheless, we also note that \textsf{FeatLLM} remains susceptible to learning spurious associations and misconstruing the causal relationships between features and the task when the attached columns come from other semantically related datasets (see Appendix~\ref{sec:spurious-appendix} for further discussion).

\vspace*{-0.12in}
\paragraph{Hyper-parameter analysis}
\textsf{FeatLLM} includes two primary hyper-parameters: the count of ensemble members and the number of rules per class extracted per LLM invocation. We systematically explored the sensitivity of model performance to these hyper-parameters. Increasing the number of ensemble components generally enhanced performance; however, the gains plateaued after approximately 20 ensembles (refer to Fig.~\ref{fig:appendix-hyperparameter1} in Appendix). Regarding the count of extracted rules, although variations did not yield statistically significant changes, our analysis suggested that setting this value to 10 provides the optimal balance between prompt complexity and predictive effectiveness (see Fig.~\ref{fig:hyperparameter2}). We conjecture that introducing more rules increases input dimensionality, which supplies additional information but also elevates the risk of overfitting, especially in low-shot scenarios.

\section{Conclusion}
We introduce \textsf{FeatLLM}, which establishes new benchmarks for LLM-driven few-shot learning on tabular data. Unlike existing strategies where LLMs directly generate predictions for each instance, \textsf{FeatLLM} emphasizes the creation of predictive criteria, aligning its design with feature engineering methodologies. Through integrating rule extraction and ensemble methods utilizing bagging, \textsf{FeatLLM} operates without heavy computational demands at inference, depends solely on API access rather than full-scale model training, and is agnostic to feature dimensionality limits. Empirical findings demonstrate that \textsf{FeatLLM} achieves competitive and often superior performance, positioning it as an efficient and practical solution for real-world tabular tasks where data resources are limited.

Currently, \textsf{FeatLLM} is optimized for low-shot learning contexts. Looking ahead, we aim to extend its applicability to more extensive datasets, develop methodologies to generate features surpassing simple rules, and introduce mechanisms to enhance interpretability. These advancements would support broader adoption across diverse use cases.

\section{Broader Impact}
The proposed framework, \textsf{FeatLLM}, seeks to harness the domain knowledge embedded in LLMs for enhanced tabular task performance, surpassing what conventional supervised methods can achieve. By leveraging this prior knowledge, the approach advances data-efficient learning and contributes strategies to reduce overfitting risks that emerge with limited labeled data. \textsf{FeatLLM} holds particular value for domains such as finance and healthcare, where acquiring labels incurs substantial costs or logistical challenges and where LLMs’ prior knowledge, trained on relevant domain texts, can offer substantive benefit. An important direction for future research involves critically examining the societal implications, including potential biases and inaccuracies present in LLM-derived knowledge, which may substantially influence prediction outputs—potentially even outweighing signals present in the available labeled samples.

\end{document}